% Part: proof-theory
% Chapter: natural-deduction
% Section: introduction

\documentclass[../../../include/open-logic-section]{subfiles}

\begin{document}

\olfileid{pt}{nat}{int}

\olsection{مقدمة}

الاستنباط الطبيعي عائلة من أنظمة~!!{proof}، لكل رابط منطقي أو سور فيها
استدلالان: أحدهما «يُدخل» العامل، وهو قاعدة~!!a{introduction}، والآخر
«يحذف» العامل، وهو قاعدة~!!a{elimination}. فمثلًا تكون نتيجة
$\Intro{\lif}$~!!{formula}s من الصورة~$!A \lif !B$، وتكون مقدمات
$\Elim{\lif}$~!!{formula}s من الصورة~$!A \lif !B$.

قدم غيرهارد غنتسن وستانيسواف ياشكوفسكي أول نسختين من الاستنباط الطبيعي
في ثلاثينيات القرن العشرين. ونظام غنتسن هو الذي يناقشه منظّرو البرهان
غالبًا، أما نظام ياشكوفسكي فكان أصل الأنظمة الأوسع استخدامًا في
التدريس. ويمكن في النظامين البدء بـ\emph{فروض}، وهي~!!{formula}s لا
يلزم أن~!!{prove}، لكن يمكن استخدامها في~!!{proving}~!!{formula}s
أخرى. وقد يعتمد~!!{proof} كامل على فروض كهذه؛ ولأنها غير مبرهنة، فلا
يثبت~!!{proof} نتيجته، أي~\emph{!!{formula} نهايته}، بل يثبت فقط أن
النتيجة تلزم من الفروض.

لبعض قواعد الاستدلال نتائج لا تعود معتمدة على بعض الفروض؛ فهي
«تبرئ» الفروض. فقاعدة~$\Intro{\lif}$ مثلًا تأخذ~!!a{proof} للنتيجة
$!B$ من الفرض~$!A$، وتكون نتيجتها~$!A \lif !B$. ولا يعود~!!{proof}
المنتهي بـ$!A \lif !B$ معتمدًا على~$!A$؛ فقد بُرئ الفرض~$!A$. وفي
أنظمة غنتسن تكون~!!{proof}s أشجارًا من~!!{formula}s، وتكون الفروض أعلى
!!{formula}s، وتحمل قواعد مثل~$\Intro{\lif}$ علامات تدل على الفروض
المبرأة. أما في أنظمة ياشكوفسكي فتُفصل، بطريقة ما، أجزاء~!!{proof}
المعتمدة على فرض؛ فقد تُزاح خلف خط رأسي أو توضع في صندوق، ويكون الفرض
$!A$ المراد تبرئته وتالي الشرط~$!B$ أول سطر في الصندوق وآخره، بينما تقع
نتيجة~$\Intro{\lif}$، أي~$!A \lif !B$، خارج الصندوق.

وتُسمى الأنظمة «طبيعية» لأن قواعد الاستدلال تحاكي الطريقة التي نستدل
بها طبيعيًا، أو التي يستدل بها الرياضيون على الأقل. فلإثبات نتيجة
افتراضية، أي شرط، نفترض المقدّم ونستخدمه لبرهنة التالي. وهذه هي
$\Intro{\lif}$. وإذا عرفنا الشرط~$!A \lif !B$ ومقدمه~$!A$، استنتجنا
$!B$. وهذه هي~$\Elim{\lif}$. ونستخدم البرهان بالحالات لنبين أن نتيجة
ما تلزم بافتراض الفصل~$!A \lor !B$. وهذه هي~$\Elim{\lor}$. ولبرهنة
دعوى عامة~$\lforall[x][!A(x)]$ نبرهن~$!A(c)$ لثابت اعتباطي~$c$. وهذه
هي~$\Intro{\lforall}$. وهلم جرًا.

وهذا مثلًا~!!{proof} بسيط لـ$!A \lif (!A \lor !B)$ في الاستنباط
الطبيعي على طريقة غنتسن:
\[
\AxiomC{$\Discharge{!A}{x}$}
\RightLabel{\Intro{\lor}}
\UnaryInfC{$!A \lor !B$}
\DischargeRule{\Intro{\lif}}{x}
\UnaryInfC{$!A \lif (!A \lor !B)$}
\DisplayProof
\]
يبدأ بالفرض~$!A$، ويستخدم~$\Intro{\lor}$ لاستنتاج~$!A \lor !B$، ثم
$\Intro{\lif}$ لاستنتاج~$!A \lif (!A \lor !B)$. ويبرئ استدلال
$\Intro{\lif}$ الفرض~$!A$، وهو ما تدل عليه العلامة~$x$.

يفضل منظّرو البرهان أنظمة غنتسن الأصلية التي تعمل بأشجار من
!!{formula}s، مع أن لأحد تنويعاتها أهمية في نظرية البرهان أيضًا.
يبين~!!^a{proof} لـ$!A$ من الفروض~$\Gamma$ أن~$!A$ تلزم من~$\Gamma$،
أي~$\Gamma \Entails !A$. ويمكن قراءة قواعد الاستدلال الطبيعي قراءة
طبيعية بوصفها تخبرنا بشيء عن علاقات اللزوم هذه؛ فمثلًا تقول
$\Intro{\lif}$ إنه إذا كان~$\Gamma + !A \Entails !B$، كان
$\Gamma \Entails !A \lif !B$. ومن الطبيعي أن نصوغ القواعد مباشرة
بحيث تعمل على الفروض والنتيجة اللازمة منها معًا في مقدمات قاعدة
الاستدلال ونتيجتها، أي أن تعمل على المتتاليات~$\Gamma \Sequent !A$.
ويأخذ البرهان البسيط السابق في نظام كهذا الصورة الآتية:
\[
\Axiom$x:!A \fCenter !A$
\RightLabel{\Intro{\lor}}
\UnaryInf$x:!A \fCenter !A \lor !B$
\DischargeRule{\Intro{\lif}}{x}
\UnaryInf$\fCenter !A \lif (!A \lor !B)$
\DisplayProof
\]
وكما ترى، القواعد هي نفسها: فهي تعمل على~!!{formula} الواقعة إلى يمين
$\Sequent$، أما~!!{formula}s الواقعة إلى اليسار فلا تفعل إلا تعداد
الفروض التي تعتمد عليها في كل خطوة.

ولا تكون المجموعة الكاملة من أزواج قواعد~!!{introduction} و
!!{elimination} تامة للمنطق الكلاسيكي بمفردها. فعلينا إضافة قاعدتين
تتعلقان بـ$\lfalse$. الأولى «قاعدة المحال الحدسية»~$\FalseInt$، أو
«الانفجار»، وتقول إننا نستطيع استنتاج أي شيء من~$\lfalse$. والثانية
مبدأ البرهان غير المباشر الكلاسيكي~$\FalseCl{}$، أو «قاعدة المحال
الكلاسيكية»، وتقول إننا إذا استطعنا أن~!!{prove}~$\lfalse$ من
$\lnot !A$، فنستطيع عندئذ أن~!!{prove}~$!A$. وإذا حذفنا~$\FalseCl$ حصلنا
على نظام
مناسب للمنطق الحدسي، وإذا حذفنا القاعدتين حصلنا على ما يسمى «المنطق
الأدنى».

سنناقش الاستنباط الطبيعي على طريقة غنتسن للمنطقين الكلاسيكي والحدسي،
أي النظامين~\Log{N1c} و\Log{N1i}. أما~\Log{N2c} و\Log{N2i} فهما
النظامان المقابلان اللذان يعملان على المتتاليات~$\Gamma \Sequent !A$
بدلًا من~!!{formula}s المفردة~$!A$.

\end{document}
